\chapter{文件系统}

\section{文件基本概念}
\concept{文件系统的组成}
\begin{points}
  \item 文件系统是操作系统中与管理文件有关的软件和数据的集合。
  \item 组成包括：被管理的文件集合、管理文件所需的数据结构、相应管理软件。
  \item 管理软件负责外存空间分配回收、目录管理、文件读写、共享与保护等。
\end{points}
\plain{文件系统这一章的主线是把“磁盘上的原始块”包装成“用户按名字访问的文件”，中间靠目录、FCB、分配方式、保护机制来衔接。}
\exam{简答题常问文件系统由什么组成，回答时最好对应“文件集合 + 数据结构 + 管理软件”三层。}


\concept{文件定义}
\begin{points}
  \item 文件是记录在外部存储上，具有名字，在逻辑上有完整意义的信息项序列。
  \item 信息项是构成文件内容的基本单位，可以是字节或记录。
  \item 从用户角度看，文件是逻辑外存的最小分配单元。
  \item 文件内容的意义由文件创建者和使用者解释。
\end{points}
\plain{用户看到的是“一个有名字、有内容的逻辑对象”，而不是磁盘块号；文件系统正是负责做这层抽象。}
\exam{名词解释题核心抓“外存上、有名字、逻辑完整、信息项序列”这几个关键词。}


\concept{文件属性}
\begin{points}
  \item 名称：便于用户识别的符号文件名。
  \item 标识符：系统内部唯一标识，常为数字形式。
  \item 类型：普通文件、目录文件、设备文件等。
  \item 位置：文件在设备上的位置指针。
  \item 大小：当前大小和最大允许大小。
  \item 保护：访问控制信息，规定谁能读、写、执行。
  \item 时间、日期、用户标识：创建、修改、访问等元数据。
\end{points}
\plain{这些属性本质上就是文件的元数据，大多保存在目录项/FCB 里，而不是文件正文里。}
\exam{选择题容易混淆“文件内容”和“文件属性”，例如把文件长度、权限、时间戳错当成正文数据。}


\concept{文件分类}
\begin{points}
  \item 按用途可分为系统文件、库文件、用户文件。
  \item 按保护级别可分为只读文件、读写文件、执行文件、不受保护文件。
  \item 按信息流向可分为输入文件、输出文件、输入输出文件。
  \item 按数据形式可分为源文件、目标文件、可执行文件。
\end{points}
\plain{分类题的关键是先看清题目按哪一个维度在分，不要把几种分类标准混成一类。}
\exam{常考选择题：问某文件属于哪一类，或者让你辨别“用途分类”和“保护级别分类”。}


\concept{文件操作}
\begin{points}
  \item Create：创建文件，为文件分配目录项和必要空间。
  \item Write：根据写指针向文件写入数据。
  \item Read：根据读指针从文件读取数据。
  \item Seek：移动文件当前读写位置。
  \item Delete：删除文件，释放空间，删除目录项。
  \item Truncate：删除文件内容但保留属性。
  \item Open/Close：打开文件时把文件控制信息装入内存打开文件表，关闭时释放相关表项。
\end{points}
\plain{`open` 的关键作用不是“开始读文件”这五个字，而是把后续读写反复要用到的控制信息先搬进内存，避免每次都重查磁盘目录。}
\exam{问 `open` 为什么重要时，重点答“减少目录搜索和磁盘访问，建立打开文件表项，保存读写状态”。}


\concept{文件类型}
\begin{points}
  \item 操作系统需要识别并支持不同文件类型，以便按合适方式处理文件。
  \item 常见类型包括普通文件、目录文件、可执行文件、设备文件、链接文件等。
  \item 文件类型信息既可能来自扩展名，也可能来自元数据或文件头。
\end{points}
\plain{文件“类型”不是纯粹给人看的标签，它会影响系统如何解释、打开和保护这个文件。}
\exam{PPT 单独列了文件类型，选择题里常把“分类标准”和“系统识别的文件类型”混在一起考。}


\section{文件访问方法与逻辑结构}
\concept{顺序访问}
\begin{points}
  \item 顺序访问按文件中信息项的先后顺序读写。
  \item 适合文本流、日志、磁带等场景。
  \item 若要访问第 $n$ 个记录，通常需要从前面依次读到它。
\end{points}
\plain{顺序访问的特点是“当前位置依赖前一步”，适合从头到尾处理文件，但不适合频繁跳着找。}
\exam{PPT 里有 \code{read\_next()}、\code{write\_next()}、\code{reset()} 这一类操作，简答时可以用来说明顺序访问接口。}


\concept{直接访问/随机访问}
\begin{points}
  \item 直接访问允许按记录号或块号直接定位到某个位置。
  \item 磁盘支持随机访问，因此文件系统可实现直接访问。
  \item 适合数据库、索引文件、需要频繁定位的数据文件。
\end{points}
\plain{直接访问的关键是能根据编号或偏移快速算出位置，所以定长记录文件尤其适合这种方式。}
\exam{PPT 特别强调了“定长记录便于直接访问，变长记录不便于直接访问”，这个点容易出判断题。}


\concept{索引访问}
\begin{points}
  \item 索引访问本质上建立在直接访问之上，通过索引表把关键字映射到记录号或偏移位置。
  \item 它适合按关键字快速查找记录，如数据库索引。
  \item 对变长记录文件，也可借助索引实现较高效定位。
\end{points}
\plain{索引访问不是顺着文件一条条扫，而是先查“目录/索引”，再跳到目标位置。}
\exam{若题目给出“关键字 -> 记录位置”的描述，通常考的就是索引访问而不是普通直接访问。}


\concept{文件逻辑结构}
\begin{points}
  \item 无结构文件/流式文件：文件被看作字节序列，如普通文本和二进制文件。
  \item 有结构文件/记录式文件：文件由一组记录组成，每条记录有内部字段结构。
  \item 顺序文件：记录按某种顺序排列，适合顺序处理。
  \item 索引文件：建立索引表，从键值映射到记录位置，提高随机查找效率。
  \item 索引顺序文件：主文件按键有序，同时建立分组索引，兼顾顺序处理和快速定位。
  \item 散列文件：通过散列函数直接计算记录位置，适合等值查找。
\end{points}
\plain{逻辑结构回答的是“用户看来文件内部怎么组织”；它和磁盘上具体放在哪些块里是两回事。}
\exam{题目若问“流式文件和记录式文件区别”，不要答成连续分配/链接分配，那已经变成物理结构了。}


\concept{文件的物理结构}
\begin{points}
  \item 文件物理结构又称存储结构，描述文件逻辑块如何映射到磁盘物理块。
  \item 它与存储设备特性和外存分配方式密切相关。
  \item 常见物理结构有顺序结构、链接结构、索引结构。
\end{points}
\plain{逻辑结构是“给用户看”的，物理结构是“给操作系统和磁盘打交道”的。}
\exam{文件逻辑结构和物理结构是经典对比点，答题时一定分清“与用户观点有关”还是“与外存组织有关”。}


\concept{顺序结构、链接结构与索引结构}
\begin{points}
  \item 顺序结构/连续结构：文件逻辑上相邻的块在磁盘上也连续，顺序访问快，但容易产生碎片，不利于动态扩充。
  \item 链接结构/串联结构：每块指向下一块，分配灵活、便于增长，但通常只适合顺序访问，且指针占空间。
  \item 索引结构：为文件建立索引表记录逻辑块与物理块映射，既支持顺序访问也支持随机访问，但有额外索引开销。
\end{points}
\plain{这三种结构本质上是在“访问效率、扩展灵活性、空间开销”三者之间做取舍。}
\exam{判断题喜欢把“链接结构便于随机访问”或“顺序结构不会产生碎片”写成陷阱。}


\section{目录结构}
\concept{磁盘结构的几个常识}
\begin{points}
  \item 磁盘通常按\key{盘块/物理块}为单位进行信息传输和空间分配，文件最终也要落实到这些磁盘块上。
  \item 一个磁盘块内部可以存放若干个逻辑记录；反过来，一个较大的逻辑记录也可能跨越多个物理块。
  \item 目录、FCB/inode、文件数据块等都只是磁盘上不同用途的数据块，只是内容和作用不同。
  \item 文件系统讨论“目录结构、物理结构、空闲块管理”时，默认背景就是这些数据最终都存放在外存块中。
\end{points}
\plain{这块不是重复大容量存储系统那章的磁头细节，而是先把文件系统这一章默认使用的“块”概念讲清：文件、目录、索引项最后都要落到磁盘块上。}
\exam{PPT 在目录结构前单独补了这层背景，选择题或简答题里若问“文件分配和传输的基本单位是什么”，通常就落在盘块/物理块。}

\concept{目录与目录项}
\begin{points}
  \item 目录用于组织系统中的文件并提供文件信息。
  \item 目录项通常包含文件名和文件控制块指针或 inode 号。
  \item 目录检索就是根据路径名逐级查找目录项，找到目标文件控制信息。
\end{points}
\plain{目录就是文件名到文件控制信息的映射表，路径就是沿目录树找到文件。}
\exam{常考一级/两级/树形目录优缺点、平均查找磁盘访问次数。}


\concept{目录管理的基本操作}
\begin{points}
  \item 目录上的基本操作通常包括：搜索文件、创建文件、删除文件、查看目录、重命名文件、遍历文件系统。
  \item 这些操作本质上都是对“文件名到目录项映射表”的查找、插入、删除或修改。
  \item 树形目录下删除目录时还要考虑删除策略：有的系统要求目录非空时不能删除，有的系统允许连同子目录和文件一起删除。
\end{points}
\plain{目录管理题不要只会说“目录能存文件名”，还要知道系统实际要支持查找、增删、改名、遍历这些操作，以及删除目录时的规则差异。}
\exam{PPT 在目录结构前单独列了目录基本操作和删除策略，选择题常考“非空目录能不能删”以及重命名/遍历属于哪类目录操作。}


\concept{分区、卷、FCB 与目录文件}
\begin{points}
  \item 分区是对物理硬盘空间的切分；原始分区需格式化生成文件系统后才能正常存放文件。
  \item 卷是已经安装好文件系统、可供读写的逻辑单元；常见场景下一个分区对应一个卷。
  \item 文件控制块 FCB 是保存文件属性和物理位置等信息的数据结构，文件与 FCB 一一对应，是文件存在的标志。
  \item 文件目录本质上是 FCB 的集合；目录文件则是把这些目录信息也作为文件保存在磁盘上。
\end{points}
\plain{这几个概念常连着考：分区/卷说明文件系统“长在什么上面”，FCB/目录说明文件系统“怎么记住每个文件”。}
\exam{PPT 对 FCB 和目录文件讲得很实，简答时可直接写“文件 = 文件体 + 文件说明（FCB）”。}


\concept{一级目录}
\begin{points}
  \item 系统只有一个目录，所有文件都放在其中。
  \item 优点：实现简单。
  \item 缺点：文件多时查找慢，不能解决重名问题，不便于多用户管理。
\end{points}
\plain{单级目录的最大问题不是“难看”，而是所有用户、所有文件都挤在一个表里，规模一大就很难管。}
\exam{答缺点时至少写出“查找慢、不能重名、不利于多用户”。}


\concept{两级目录}
\begin{points}
  \item 每个用户有自己的用户文件目录 UFD，系统有主文件目录 MFD。
  \item 解决不同用户之间文件重名问题。
  \item 用户内部文件多时仍不够灵活。
\end{points}
\plain{二级目录解决的是“不同用户同名文件冲突”，但对同一用户内部的大量分类管理帮助仍有限。}
\exam{PPT 还提到它不利于合作共享，这也是两级目录常见缺点。}


\concept{树形目录与路径}
\begin{points}
  \item 树形目录把目录组织成层次结构，便于分类管理。
  \item 绝对路径从根目录开始，如 \code{/home/a.txt}。
  \item 相对路径从当前目录开始。
  \item 优点：层次清晰、重名控制灵活、便于权限管理。
\end{points}
\plain{树形目录的关键是“路径”与“当前目录”两个概念：同一个文件既能用绝对路径定位，也能结合当前目录用相对路径定位。}
\exam{PPT 还提了特殊目录 `.` 和 `..`，以及绝对/相对路径的区别，这些很适合出选择题。}

\concept{当前目录与特殊目录项}
\begin{points}
  \item \term{当前目录}（working directory）是用户在某一段时间内默认工作的目录。
  \item 相对路径不是从根目录出发，而是从当前目录出发解释。
  \item 特殊目录项 \code{.} 表示当前目录本身，\code{..} 表示当前目录的父目录。
  \item 这两个特殊目录项使用户能在目录树中方便地进行相对定位与层次切换。
\end{points}
\plain{当前目录可以理解成“你现在站在哪个目录里”；有了它，很多路径就不必每次都从根目录写全。}
\exam{选择题常直接问 \code{.} 和 \code{..} 各表示什么，或问相对路径是相对于根目录还是相对于当前目录。}


\concept{无环图目录与共享}
\begin{points}
  \item 无环图目录允许多个目录项共享同一个文件或子目录，但不允许形成环。
  \item 共享文件需要引用计数，记录有多少目录项指向该文件。
  \item 引用计数为 0 时，文件数据和控制信息才能真正回收。
\end{points}
\plain{无环图目录是在树形目录基础上加入“共享同一个文件/目录”的能力，但又必须防止形成环导致路径遍历死循环。}
\exam{容易考“为什么要引用计数”“为什么不允许形成环”，不要答成并发互斥题。}


\concept{通用图目录}
\begin{points}
  \item 若在目录树中不断增加链接，可能形成一般图结构。
  \item 通用图目录允许更灵活的共享，但会带来循环搜索、删除回收更复杂等问题。
  \item 为避免无限遍历，系统通常需要限制搜索深度或在遍历时记录已访问结点。
\end{points}
\plain{图结构比无环图更自由，但自由的代价是目录遍历、删除和垃圾回收都更麻烦。}
\exam{PPT 单独点出它的风险是“循环搜索”和“删除复杂”，这是判断题常见考点。}


\section{文件系统安装、共享与保护}
\concept{文件系统安装 mount}
\begin{points}
  \item 文件系统必须挂载到某个目录位置后，用户才能通过统一目录树访问它。
  \item 挂载点是原有目录树中的一个目录。
  \item 挂载后，该目录下呈现被挂载文件系统的根内容。
\end{points}
\plain{挂载的本质是把“另一个卷上的文件系统根目录”接到当前目录树的某个节点下面。}
\exam{若题目给出 `mount` 场景，关键点是“统一命名空间”和“挂载点目录”。}


\concept{多用户共享中的所有者与组}
\begin{points}
  \item 多用户系统通常用所有者 owner 和组 group 来组织文件共享与保护。
  \item 所有者通常拥有修改文件属性和授予访问权限的最高控制权。
  \item 组表示一组用户，系统可按组统一授予读、写、执行等权限。
  \item 文件或目录会保存用户标识 UserID 和组标识 GroupID，访问时据此判断权限。
\end{points}
\plain{这套机制的核心是“先判你是不是文件主人，再判你是不是这个组的人”，不同身份走不同权限规则。}
\exam{PPT 在文件共享里专门先讲 owner/group，选择题常考 UserID、GroupID 和组权限分别起什么作用。}


\concept{硬链接与符号链接}
\begin{points}
  \item 硬链接是多个目录项指向同一个 inode/文件控制块，本质上是同一个文件有多个名字。
  \item 删除一个硬链接只减少引用计数；只有引用计数为 0 时文件才真正删除。
  \item 符号链接是一个特殊文件，内容是目标文件路径。
  \item 符号链接可以跨文件系统，也可以指向目录；但目标删除后可能变成悬空链接。
\end{points}
\plain{硬链接共享同一个 inode，符号链接是另一个文件，内容是目标路径。打开文件表保存运行时读写状态。}
\exam{常考硬/软链接区别、引用计数、系统级/进程级打开文件表。}


\concept{共享一致性语义}
\begin{points}
  \item 一致性语义描述多个用户同时访问共享文件时，一个用户的修改何时对其他用户可见。
  \item Unix 语义通常认为写操作对其他已打开同一文件的用户立即可见。
  \item AFS 会话语义通常在文件关闭后，修改才对之后新打开该文件的会话可见。
  \item 不可变共享文件语义则要求共享文件一旦发布后内容不可再修改。
\end{points}
\plain{共享不只是“能不能一起用”，还包括“我改完之后别人什么时候看到我的改动”。}
\exam{这类题容易被忽视，但 PPT 单独列了共享一致性语义，至少要会区分 Unix 立即可见和 AFS 会话可见。}


\concept{远程文件系统与故障模式}
\begin{points}
  \item 远程文件访问常见有三类方式：手工程序传输方式（如 FTP）、自动化无缝访问方式（如 NFS、SMB、HDFS 一类 DFS）、以及基于 Web 的半自动化访问方式。
  \item DFS 的特点是可把远端存储挂载到本地目录树中，程序可像访问本地文件那样透明访问远程文件。
  \item 纯 Web/HTTP 方式通常更适合获取和展示文件，不天然提供透明随机写；但配合 WebDAV 等扩展后也可接近挂载访问。
  \item 与本地文件系统相比，远程文件系统除硬件故障、元数据损坏、数据损坏外，还会增加网络故障和服务器故障。
  \item 故障恢复时常涉及请求状态管理，因此会出现无状态 DFS 与有状态 DFS 的区分。
\end{points}
\plain{远程文件系统把“磁盘问题”扩展成了“磁盘 + 网络 + 服务器”三类问题，所以故障模式更复杂。}
\exam{老师录音没展开，但 PPT 有这一块；若出选择题，常考远程文件系统新增的故障来源。}


\concept{文件锁}
\begin{points}
  \item 文件锁允许进程锁定文件，以防其他进程在不合适的时机访问它。
  \item 共享锁类似读锁，允许多个进程并发读文件。
  \item 独占锁类似写锁，同一时刻只允许一个进程持有。
  \item 强制锁由操作系统直接阻止冲突访问；建议锁只提供机制，要求程序自行遵守。
\end{points}
\plain{文件锁本质上就是把读者写者问题搬到了文件访问上。}
\exam{常考共享锁和独占锁的区别，或问强制锁/建议锁到底由谁保证生效。}


\concept{文件保护}
\begin{points}
  \item 文件保护用于控制用户对文件的读、写、执行、删除等权限。
  \item 文件保护既要防止系统崩溃带来的破坏，也要防止用户有意或无意的非法操作。
  \item 常见方法包括访问控制表 ACL、权限位、口令保护、能力表等。
  \item Unix 风格权限常分为所有者、同组用户、其他用户三类，每类有读写执行位。
\end{points}
\plain{文件保护一部分是“防人乱动”，另一部分是“防系统崩了把文件弄坏”；前者偏权限控制，后者偏备份容错。}
\exam{PPT 把这两类问题分开讲了，简答题不要只写访问控制，还可补“定时转储、多副本”等保护思路。}


\concept{文件保护的两类问题}
\begin{points}
  \item 第一类问题是防止系统崩溃或介质故障导致文件被破坏。
  \item 针对这类问题，常见办法有定时转储、建立多副本、备份关键目录和分配表等。
  \item 第二类问题是防止文件所有者或其他用户有意、无意地做出非法操作。
  \item 针对这类问题，核心机制是访问控制，即按“用户 - 对象 - 权限”决定能否访问。
\end{points}
\plain{文件保护题别只盯着权限位；PPT 的原始分法是先问“防系统坏”还是“防人乱动”，然后再分别给办法。}
\exam{简答题若问“文件保护包括哪些内容”，最好先按这两大类分开答，再补 ACL、权限位、备份、多副本等具体手段。}


\concept{打开文件表}
\begin{points}
  \item 打开文件后，系统把文件控制信息缓存在内存打开文件表中，避免每次读写都重新搜索目录。
  \item 系统级打开文件表记录文件位置、引用计数、磁盘位置等全局信息。
  \item 进程级打开文件表记录进程自己的文件描述符、访问模式、当前读写指针等。
  \item 若父子进程共享同一个打开文件表项，可能共享文件偏移；若分别打开同一文件，则偏移独立。
\end{points}
\plain{打开文件表可以理解成“把磁盘上的静态目录信息转成内存里的运行时状态”，这样后续读写才高效。}
\exam{PPT 明确讲了“局部表 + 全局表”，要会区分哪些信息按进程独有，哪些信息全系统共享。}
